\section[Metrized Cubic Algebras]{Metrized Cubic Algebras and Differential Closure}
\label{sec:cubic-algebraic-theory}

For an essential irreducible harmonic cubic with minimal regular zero
set, the smallest Laplacian algebra has at most four generators in the
radial case and exactly five in the nonradial case
(\cref{thm:harmonic-cubic-five-bound}). Radial transport identities
and the Clifford table prepare \cref{sec:cubic-potentials}.
The nonradial weight operator gives an intrinsic graded splitting,
five differential invariants, and a reduction to Cartan cubics in
dimensions $8,14,26$ (\cref{cor:nr-cartan-dimensions}); compatible
splittings still require classification. For the five-generator proof,
pass from \cref{subsec:cubic-multiplication} to
\cref{sec:nonradial-cubics,subsec:nr-recovery,subsec:nr-five-generators};
coordinate models appear in \cref{sec:nonradial-examples}.

\subsection[Hsiang's problem and radial cubic classification]
{Hsiang's radial cubic problem and the classification work of
Peng--Xiao and Tkachev}
\label{sec:cubic-history}

Within his broader study of algebraic minimal cones, Hsiang
\cite{Hsiang67} singled out the classification of cubic forms satisfying
\[
 \cL(F)=-\gamma RF,\qquad R=|x|^2,
\]
with a constant $\gamma$. These are now called \emph{radial
eigencubics}; radiality refers to the quadratic weight in this equation,
not to radial dependence of $F$. This is the radial cubic classification
problem considered here; see \cite[Section~1.1]{TkachevClassification10}
and \cite[Section~1]{FoxTkachev24}.
Hsiang's invariant-theoretic constructions recover Cartan cubics in
$\R^5$ and $\R^8$ and give the $SO(4)$- and $SU(4)$-invariant
models in $\R^9$ and $\R^{15}$; see
\cite[Section~6.1]{TkachevClassification10}.
Tongyan Liu \cite{Liu89} later gave an explicit octonionic construction
of a further cubic in $\R^{27}$, the $Sp(4)$-invariant model;
see \cite[Example~4 and Theorem~C]{PengXiao93} and
\cite[Section~6, p.~1562]{HoppeTkachev19}.

Peng--Xiao \cite{PengXiao93} and, independently, Tkachev
\cite{TkachevClifford10,TkachevClassification10} developed closely
related approaches to radial minimal cubics.\footnote{Tkachev records
that he learned of Peng--Xiao's 1993 paper from Zizhou Tang in November
2010; see \cite[Remark~1, p.~40]{TkachevMiklyukov19}.}
Peng--Xiao assumed harmonicity; Tkachev proved it for irreducible radial
cubics \cite[Theorem~1]{TkachevClassification10}; the algebraic
form of this result, with cubes of linear forms as the exception, is
\cite[Theorem~6.6.1]{NadirashviliTkachevVladut14}.

\paragraph{Clifford and exceptional cubics.}
Following Tkachev, an irreducible radial cubic is of \emph{Clifford
type} if it is orthogonally equivalent, up to a nonzero constant factor,
to
\[
 F_{\mathcal A}(y,z)=\sum_{i=0}^{q}z_i\langle A_i y,y\rangle,
 \qquad A_iA_j+A_jA_i=2\delta_{ij}\Id,
\]
where $y\in\R^{2m}$, $z\in\R^{q+1}$, and the $A_i$ are symmetric
matrices on $\R^{2m}$. The remaining irreducible
radial cubics are called \emph{exceptional}; see
\cite[Section~1.2]{TkachevClassification10} and
\cite[Definitions~2.2 and~2.8]{FoxTkachev24}.
In particular, Cartan cubics belong to the exceptional branch.
The Clifford family will be written explicitly in
\cref{sec:invariants} to compute its differential closure.

% Included in the historical opening subsection of Section 9.
\paragraph{Radial normal forms and the classical model families.}
\label{rem:cubic-radial-normal-form-history}
Choosing a positive maximum direction on the unit sphere, with
coordinate $z$, and normalizing $F$ gives
\[
 \begin{gathered}
 F=z^3+zB+C,\qquad
 B=-3|\xi|^2-\tfrac32|\eta|^2+\tfrac32|\zeta|^2,\\
 C=C_{111}+C_{102}+C_{012}+\sigma,\qquad \sigma=C_{030}.
 \end{gathered}
\]
The subscripts denote multidegrees in $(\xi,\eta,\zeta)$, whose
dimensions are $n_1,n_2,n_3$, respectively. They satisfy
$n_3=2n_1+n_2-2$ and
$N=3n_1+2n_2-1$; see \cite[Section~1, Theorem~1]{PengXiao93}
and \cite[Proposition~3.1]{TkachevClassification10}.

The component $\sigma$ satisfies
$|\nabla\sigma|^2=\tfrac92|\eta|^4$.
The zero case means that the $\eta$-space is absent ($n_2=0$).
Otherwise Peng--Xiao \cite[Theorem~2]{PengXiao93} classify $\sigma$
as reducible or a scaled Cartan cubic; a different proof of the
eikonal classification is given in
\cite[Theorem~1.1]{TkachevCartan10}.
The case $\sigma=0$ yields the Cartan cubics, whereas the case
$\sigma\ne0$ with $\sigma$ reducible yields the Clifford cubics
\cite[Theorems~A and~B]{PengXiao93}; see also
\cite[Proposition~3.4 and Theorem~3]{TkachevClassification10}.
In the remaining branch, $n_2\in\{5,8,14,26\}$.
The cases $n_1=0$ and $n_1=1$ were treated in
\cite[Theorems~C and~D]{PengXiao93} and
\cite[Sections~6.2--6.3]{TkachevClassification10}:
the former yields the three models invariant under $SO(4)$, $SU(4)$,
and $Sp(4)$, and the latter yields four models interpreted through
complexified rank-three Jordan norms
\cite[Section~6.11.6(iii)--(iv)]{NadirashviliTkachevVladut14}.

Tkachev also identifies a hidden Clifford system encoded by
$C_{111}+3^{-1/2}C_{102}$ when $n_1>0$, imposing Hurwitz--Radon restrictions on the
multiplicities \cite[Proposition~3.3]{TkachevClassification10}.
His construction of a $G_2$-invariant cubic of type $(4,5)$ and his
correspondence with isoparametric quartics provide further examples and
restrictions \cite[Section~6.4 and Theorem~5]{TkachevClassification10},
using the isoparametric classification results of
Cecil--Chi--Jensen \cite{CecilChiJensen07} and Immervoll
\cite{Immervoll08}. The exceptional branch has only finitely many
possible multiplicity types, with $N\le72$
\cite[Theorem~6.11.2]{NadirashviliTkachevVladut14}.
This bounds the possible types rather than completing the
classification of their realizations.



Subsequent work of Tkachev and his collaborators recast radial minimal
cubics in the language of commutative nonassociative metrized algebras,
giving intrinsic interpretations of their Clifford and Jordan
structures \cite{TkachevJordan14,NadirashviliTkachevVladut14,FoxTkachev24}.
We follow this terminology and recall the algebraic preliminaries
needed to study their Laplacian algebras.

\subsection{Metrized multiplication and low-degree structure}
\subsubsection{The multiplication associated with a cubic}
\label{subsec:cubic-multiplication}

The radial equation above also defines the associated
\emph{radial eigencubic (REC) algebras}
\cite[Definition~6.4.1]{NadirashviliTkachevVladut14}.
We allow an arbitrary quadratic weight in \(\cL(F)=-qF\).

We use the \emph{Freudenthal--Springer algebra} \(V(F)\) of a cubic
\cite[Definition~6.3.1 and equations~(6.4.1)--(6.4.4)]{NadirashviliTkachevVladut14}.
Let \(V\) be a Euclidean space and let \(F\in\R[V]\) be homogeneous
of degree three. Define a bilinear multiplication by
\begin{equation}\label{eq:cubic-polarized-product}
 \langle xy,z\rangle=D^3F[x,y,z],\qquad x,y,z\in V.
\end{equation}
The inner product determines \(xy\); symmetry of \(D^3F\) gives
\[
 xy=yx,\qquad \langle xy,z\rangle=\langle x,yz\rangle.
\]
Thus \(V(F)\) is commutative with invariant Euclidean inner product.
Its \emph{Freudenthal multiplication} gives self-adjoint operators
\(L_xy=xy\). No unit, associativity, or Jordan identity is assumed;
in particular, \(L_{x^2}y=(xx)y\) and \(L_x^2y=x(xy)\) need not agree.

\begin{proposition}[Differential quantities and multiplication]
\label{prop:cubic-multiplication-dictionary}
With the normalization \eqref{eq:cubic-polarized-product},
\begin{equation}\label{eq:cubic-multiplication-dictionary}
 F(x)=\frac16\langle x^2,x\rangle,\qquad
 \nabla F(x)=\frac12x^2,\qquad
 \Hess F(x)=L_x,\qquad \Delta F(x)=\tr L_x.
\end{equation}
Here the Hessian is identified with a self-adjoint endomorphism using
the Euclidean inner product. Conversely, a commutative bilinear
multiplication with an invariant Euclidean inner product determines a
cubic by the first formula, and its polarization recovers the multiplication.
\end{proposition}

\begin{proof}
Homogeneity gives \(D^3F[x,x,x]=6F(x)\). Differentiating the
first formula and using invariance of the inner product gives
\[
 DF(x)[v]=\frac16\bigl(2\langle xv,x\rangle+
                         \langle x^2,v\rangle\bigr)
        =\frac12\langle x^2,v\rangle.
\]
Another derivative gives \(D(\nabla F)(x)[v]=xv\); taking its
trace proves the last assertion. For the converse, commutativity and
invariance make \(\langle xy,z\rangle\) fully symmetric, so the
third derivative of \(\langle x^2,x\rangle/6\) is precisely this
trilinear form.
\end{proof}

Writing \(x^3=x(x^2)\), substitution in
\(\cL(F)=|\nabla F|^2\Delta F-
\Hess F(\nabla F,\nabla F)\) shows that, for a quadratic form \(q\),
the equation \(\cL(F)=-qF\) is equivalent to
\[
 |x^2|^2\tr L_x-\langle x^2,x^3\rangle
 =-\frac23q(x)\langle x^2,x\rangle.
\]
If \(F\) is harmonic, then \(\tr L_x=0\) and the identity reduces to
\(\langle x^2,x^3\rangle=\frac23q(x)\langle x^2,x\rangle\).
The further specialization \(q=\gamma R\) gives the radial REC
identity in \cite[Proposition~6.4.1]{NadirashviliTkachevVladut14},
which uses the same cubic normalization and the parameter
\(\lambda=-\gamma\).

The metrized algebra on \(V\) encodes the derivatives of \(F\), whereas
\(\cA_R(F)\) contains scalar differential invariants such as
\(S=|\nabla F|^2=|x^2|^2/4\). In particular, the critical points
of \(F\) are exactly the square-zero elements of its metrized algebra.

The cylindrical directions are encoded by
\[
 D_vF(x)=\frac12\langle x^2,v\rangle
        =\frac12\langle L_vx,x\rangle,
\]
which gives \(V_{\mathrm{cyl}}(F)=\{v:L_v=0\}\), the annihilator
of its metrized algebra.

\subsubsection{The quadratic weight and linear generators}

An irreducible cubic with nonempty minimal regular zero set satisfies,
by \cref{prop:divisibility},
\begin{equation}\label{eq:cubic-general-weight}
 \cL(F)=-qF,
\end{equation}
where \(q=-\lambda_F\) is a homogeneous quadratic polynomial,
possibly zero. Irreducibility is used here to obtain the polynomial
equation; its radial specialization was introduced in
\cref{sec:cubic-history}.

For brevity put \(\cB=\cA_R(F)\) and
\(\cB_d=\cB\cap\R[V]_d\), the degree-\(d\) homogeneous part.

\begin{proposition}[Quadratic information in the closure]
\label{prop:cubic-weight-in-minimal-algebra}
Suppose \(\dim V\geq2\) and \(F\not\equiv0\) is a homogeneous
cubic satisfying \eqref{eq:cubic-general-weight} with a polynomial
\(q\). Then \(q\in\cB_2\). If \(\cB_2=\R R\), then
\[
 \Delta F=0,\qquad q=\gamma R\quad\text{with }\gamma>0.
\]
In particular, a nonradial weight forces a quadratic element of
\(\cA_R(F)\) independent of \(R\).
\end{proposition}

\begin{proof}
The identity \(\cL(F)=S\Delta F-\tfrac12\Gamma(F,S)\)
shows that \(\cL(F)\in\cB\). The polynomial-quotient property
in \cref{prop:MR-finite-and-division} therefore gives
\(q=-\cL(F)/F\in\cB_2\). If \(\cB_2=\R R\), the linear
polynomial \(\ell=\Delta F\) satisfies \(\ell^2\in\R R\).
A nonzero square of a linear form has rank one, whereas \(R\)
has rank \(\dim V\geq2\); hence \(\ell=0\).
We already know \(q=\gamma R\). At a positive maximum \(a\) of
\(F\) on the unit sphere, attained at \(x\), homogeneity gives
\(\nabla F(x)=3ax\) and \(\Hess F(x)x=6ax\).
Evaluating \eqref{eq:cubic-general-weight} there gives
\(-\gamma a=-54a^3\), so \(\gamma=54a^2>0\).
\end{proof}

\begin{remark}[Additional quadratic generators]
\label{rem:cubic-nonradial-closure}
Since \(\R[R,F,S]_2=\R R\),
\cref{prop:cubic-weight-in-minimal-algebra} shows that its Laplacian
closure forces radiality. Adjoining a nonradial \(q\) need not close
the algebra. Conversely, the radial unbalanced Clifford cubics have
several quadratic generators; see \cref{thm:clifford-minimal-algebra}.
\end{remark}

\begin{remark}[Linear generators in the nonharmonic case]
\label{rem:cubic-nonharmonic-linear-splitting}
If \(\Delta F\ne0\), then \(\cB_1\ne0\). Choose orthonormal
linear coordinates \(t_1,\ldots,t_r\) spanning \(\cB_1\), and put
\(W=\operatorname{span}\{\nabla t_i\}\). Then
\[
 \cB=\R[t_1,\ldots,t_r]\otimes\cB',\qquad
 \cB'=\cB\cap\R[W^\perp],\qquad \cB'_1=0.
\]
Indeed, \(\partial_{t_i}h=\Gamma(t_i,h)\in\cB\), so the finite
Taylor identity
\[
 h(0,y)=\sum_\alpha\frac{(-t)^\alpha}{\alpha!}\partial_t^\alpha h(t,y),
\]
applied also to the derivatives of \(h\), puts every coefficient
in \(\cB'\). This algebra is Laplacian and contains
\(|y|^2=R-\sum_i t_i^2\). The cubic \(F\) can still have mixed
terms, and the splitting gives no uniform generator bound for \(\cB'\).
\end{remark}

\subsubsection{One cubic generator and orthogonal blocks}

The quadratic part of a Laplacian algebra is a Jordan algebra of
self-adjoint endomorphisms; see Mendes--Radeschi
\cite{MendesRadeschiQuadratic}. The following restriction on its
cubic part forces its endomorphisms to commute.

\begin{theorem}[Orthogonal blocks forced by the cubic part]
\label{thm:cubic-one-dimensional-degree-three}
Let \(F\not\equiv0\) be an essential homogeneous cubic on a Euclidean
space \(V\) of dimension at least two, and suppose
\(\cB_3=\R F\). Then \(\cB_1=0\) and \(\Delta F=0\).
There is an orthogonal decomposition into nonzero subspaces,
unique up to permutation,
\begin{equation}\label{eq:cubic-canonical-blocks}
 V=V_1\oplus\cdots\oplus V_r,\qquad
 \cB_2=\operatorname{span}\{R_1,\ldots,R_r\},\qquad
 R_i(x)=|\operatorname{pr}_{V_i}x|^2.
\end{equation}
Moreover, \(F\) is multihomogeneous of positive degrees
\((d_1,\ldots,d_r)\), with \(d_1+\cdots+d_r=3\).
Consequently \(r\leq3\), and, up to permutation, the possible degree
patterns are
\[
 (3),\qquad (2,1),\qquad (1,1,1).
\]
If \(F\) also satisfies \eqref{eq:cubic-general-weight}, then
\(q=\sum_i\gamma_iR_i\) for constants \(\gamma_i\).
\end{theorem}

\begin{proof}
If \(0\ne\ell\in\cB_1\), then both \(\ell R\) and \(\ell^3\)
are nonzero elements of the one-dimensional space \(\cB_3\).
Cancellation would give \(\ell^2=cR\), contradicting the rank
argument above. Thus \(\cB_1=0\), and in particular \(\Delta F=0\).

Represent \(\cB_2\) by a subspace \(J\) of self-adjoint
endomorphisms: \(B\in J\) when \(q_B(x)=\langle Bx,x\rangle\)
belongs to \(\cB_2\). Since
\[
 \Gamma(q_B,q_C)=4q_{B\circ C},\qquad
 B\circ C=\tfrac12(BC+CB),
\]
\(J\) contains \(\Id\) and is closed under the Jordan product.
For \(B\in J\), define \(E_Bh(x)=\langle Bx,\nabla h(x)\rangle\).
Then \(2E_BF=\Gamma(q_B,F)\in\cB_3\), so
\(E_BF=\mu_BF\) for some real number \(\mu_B\).

We claim that all matrices in \(J\) commute. Fix \(B,C\in J\)
and put \(K=[B,C]\). The identity
\([E_B,E_C]=-E_K\) gives \(E_KF=0\). Since \(K\) is
skew-adjoint, the orthogonal transformations \(\exp(tK)\) preserve
\(F\). Every orthogonal symmetry of \(F\) fixes every element of
\(\cB\): its invariant polynomial algebra is Laplacian and contains
\(R,F\), and hence contains \(\cA_R(F)\). In particular,
\(q_B(\exp(tK)x)=q_B(x)\), which upon differentiation gives
\([B,K]=0\). In an orthonormal basis diagonalizing \(B\),
the \((i,j)\) entry of \([B,[B,C]]\) is
\((b_i-b_j)^2C_{ij}\). Thus \([B,C]=0\), as claimed.

Simultaneously diagonalize \(J\), and group together the coordinates
on which every member of \(J\) has the same diagonal value. These
groups define the joint eigenspaces \(V_1,\ldots,V_r\).
A generic element of \(J\) has distinct eigenvalues on distinct
groups. Its spectral projections \(P_i\) are polynomials in this
element and therefore belong to \(J\), because on commuting
matrices the Jordan product is ordinary multiplication. Hence
\(J=\operatorname{span}\{P_1,\ldots,P_r\}\).
The projections are precisely the primitive nonzero idempotents of
this commutative matrix algebra, proving uniqueness in
\eqref{eq:cubic-canonical-blocks}.

Finally, \(E_{P_i}\) is the Euler operator in the variables of
\(V_i\), and \(E_{P_i}F=d_iF\) for a real scalar \(d_i\).
Expansion into monomials shows that each \(d_i\) is a nonnegative
integer and that \(F\) has degree \(d_i\) in that block.
Total homogeneity gives \(\sum_i d_i=3\). If \(d_i=0\), \(F\)
would be independent of \(V_i\), so
\(V_i\subset V_{\mathrm{cyl}}(F)=\{0\}\), contradicting
\(V_i\ne\{0\}\). This proves the degree list. The final assertion follows
from \cref{prop:cubic-weight-in-minimal-algebra}.
\end{proof}

Clifford cubics can also require a quartic generator
(\cref{thm:clifford-minimal-algebra}). The assumption
\(\cB_3=\R F\) is absent from the subcalibration criteria below
and fails for the nonradial models, whose two cubic components are
separated by a block Euler operator.

\subsection{Critical Hessians and multiplier degeneracy}
\label{sec:cubic-square-zero-obstruction}

By \cref{prop:cubic-multiplication-dictionary}, critical points are
square-zero elements \(c^2=0\), and \(L_c=\Hess F(c)\) governs
the first transverse term of \(F\). The Hessian identity below follows
from \eqref{eq:cubic-general-weight} for an arbitrary quadratic weight
\(q\).

\begin{proposition}[Critical Hessians and multiplier degeneracy]
\label{prop:cubic-general-singular-spectrum}
\label{prop:cubic-positive-multiplier-obstruction}
\label{cor:cubic-essential-multiplier-obstruction}
Let $F$ be a nonzero homogeneous cubic on $\R^N$ satisfying
\eqref{eq:cubic-general-weight}, where $q$ is a quadratic form. Let $c$ be a critical point, and put
$A=\Hess F(c)$ and $\tau=\tr A$. Then
\begin{enumerate}[label=\textup{(\roman*)}]
\item The Hessian satisfies
\begin{equation}\label{eq:cubic-general-hessian-polynomial}
 A^3-\tau A^2=\frac{q(c)}2 A.
\end{equation}
In particular, every nonzero eigenvalue satisfies
$\lambda^2-\tau\lambda-q(c)/2=0$. If $q(c)\ne0$, then $A\ne0$.
If $F$ is harmonic and $A\ne0$, then $q(c)>0$ and the nonzero eigenvalues are
$\pm\sqrt{q(c)/2}$, with equal multiplicities.
\item If $q(c)>0$ and $a$ is $C^2$ near $c$, with $a(c)>0$,
then $G=aF$ cannot satisfy $F\cL(G)\ge0$ throughout a neighborhood
of $c$.
\item If $F$ is essential and harmonic, every nonzero critical point
has $A\ne0$ and $q(c)>0$. Hence a nonnegative smooth multiplier
giving the two-sided subcalibration sign must vanish at every such
point across which it extends smoothly.
\end{enumerate}
\end{proposition}

\begin{proof}
For (i),
Euler's identity gives $F(c)=0$. Write
\[
 F(c+y)=h(y)+T(y),\qquad
 h(y)=\tfrac12\langle Ay,y\rangle,
\]
where $T$ is homogeneous of degree three and is not identically zero.
If $q(c)\ne0$ and $A=0$, the left side of \eqref{eq:cubic-general-weight}, expanded
at $c$, is homogeneous of degree five in $y$, whereas the right side
has the nonzero term $-q(c)T(y)$ of degree three. This is impossible.

The terms of degree two in the same equation are
\[
 \tau\langle A^2y,y\rangle-\langle A^3y,y\rangle
 =-\frac{q(c)}2\langle Ay,y\rangle.
\]
All matrices involved are symmetric, so this proves
\eqref{eq:cubic-general-hessian-polynomial}. If $F$ is harmonic,
$\tau=0$. If in addition $A\ne0$, a nonzero real eigenvalue gives
$q(c)/2=\lambda^2>0$. The zero trace then forces the multiplicities
of the two nonzero eigenvalues to agree.

For (ii), \cref{lem:critical-multiplier-defect} gives
\(\cK_F(a)(c)=-a(c)^3q(c)<0\). Since (i) gives \(A\ne0\),
choose \(y\) with \(h(y)\ne0\). Continuity and
\(F(c+ry)=r^2h(y)+O(r^3)\) yield
\begin{equation}\label{eq:cubic-positive-multiplier-leading-sign}
 F(c+ry)\cL(aF)(c+ry)
 =-a(c)^3q(c)r^4h(y)^2+o(r^4)<0
\end{equation}
for small \(r>0\). The negative product implies
\(\nabla(aF)\ne0\), so the sign fails where the field is defined;
harmonicity is not used.

For (iii), \(L_c=0\) would put \(c\) in the annihilator
\(V_{\mathrm{cyl}}(F)=\{0\}\). Thus (i) gives \(q(c)>0\).
The multiplier conclusion is also the harmonic obstruction of
\cref{prop:harmonic-positive-multiplier-obstruction}.
\end{proof}

For essential harmonic $F$ as above, \cref{sec:tangent-instability-transfer}
identifies the tangent at every nonzero critical point as
$\{(u,v)\in\R^s\oplus\R^s:|u|=|v|\}\times\R^{N-2s}$,
where $2s$ is the Hessian rank, with smooth convergence on compact
regular subsets. Radial powers $R^\beta$ fail as smooth positive
multipliers there, whereas $S$ has the required vanishing.
The following trace bound supplies removability for
\cref{thm:cubic-unified-sextic}.

\begin{lemma}[A trace bound for the critical set]
\label{lem:cubic-trace-removability}
Let \(F\not\equiv0\) be a harmonic cubic on \(\R^N\), \(N\ge2\),
satisfying \(\cL(F)=-2RF\), and put \(Z=\{\nabla F=0\}\).
If \(\Delta S\ge16R\), then every \(c\in Z\setminus\{0\}\)
satisfies \(\operatorname{rank}\Hess F(c)\ge8\), and \(Z\) is locally
contained near \(c\) in a smooth submanifold of codimension eight.
In particular, \(Z\) has locally zero \(\cH^{N-1}\)-measure.
\end{lemma}

\begin{proof}
At \(c\in Z\setminus\{0\}\),
\cref{prop:cubic-general-singular-spectrum} gives
\(A^3=|c|^2A\), where \(A=\Hess F(c)\). Its nonzero eigenvalues
are \(\pm|c|\), and Bochner's identity yields
\[
 |c|^2\operatorname{rank}A=\tr(A^2)
   =\tfrac12\Delta S(c)\ge8|c|^2.
\]
Choose eight components of \(\nabla F\) with independent differentials.
The implicit function theorem places \(Z\) locally in their smooth
common zero set of codimension eight. This proves the local codimension
claim and, by a countable cover, nullity off the origin. The origin is
also null since \(N\ge2\).
If \(N<8\), the rank inequality simply excludes nonzero critical points.
\end{proof}

\subsection{Radial transport and Clifford closure}
\subsubsection{Intrinsic transport and the scalar trace form}

For harmonic \(F\), the identity \(\Gamma(F,S)=4RF\) is the
normalization \(\lambda=-2\) for REC algebras. We use this normalization
and its polarized identities from
\cite[Section~6.7.1, equations~(6.7.6)--(6.7.7)]{NadirashviliTkachevVladut14}.

\begin{lemma}[Intrinsic radial transport identities]
\label{lem:cubic-radial-transports}
Let \(F\) be a homogeneous cubic on a Euclidean space \(V\), with
\(\Gamma(F,S)=4RF\), where \(R=|x|^2\) and \(S=|\nabla F|^2\).
Define
\begin{equation}\label{eq:cubic-unified-D}
 D=\frac{\Gamma(S,S)-4RS}{12}
   =\frac{|x^3|^2-R|x^2|^2}{12}.
\end{equation}
Then
\begin{equation}\label{eq:cubic-unified-table}
 \Gamma(F,D)=2SF,\qquad
 \Gamma(S,D)=8RD,\qquad
 \Gamma(D,D)=4SD.
\end{equation}
Neither harmonicity nor a scalar quadratic trace form is needed.
\end{lemma}

\begin{proof}
Use the multiplication \eqref{eq:cubic-polarized-product}, and put
\[
 \begin{gathered}
 y=x^2,\quad z=xy,\quad u=y^2,\quad w=yz,\\
 r=|x|^2,\quad c=\langle x,y\rangle,\quad
 s=|y|^2,\quad t=|z|^2,\quad \delta=t-rs.
 \end{gathered}
\]
Thus \(F=c/6\), \(S=s/4\), \(D=\delta/12\), and the hypothesis is
\(P(x):=\langle y,z\rangle-4rc/3=0\). Its gradient gives the
polarized identity
\begin{equation}\label{eq:cubic-radial-polarization}
 Q(x):=4xz+u-\tfrac83cx-4ry=0.
\end{equation}
Pairing with \(y\) gives
\begin{equation}\label{eq:cubic-radial-polarization-eta}
 \eta:=\langle y,u\rangle=\tfrac83c^2-4\delta.
\end{equation}
The coefficient of \(\tau^2\) in \(P(x+\tau y)=0\), with \(y\)
held fixed, and the pairing \(\langle Q,z\rangle=0\) give, respectively,
\[
 \begin{aligned}
 4\langle z,xz\rangle+6\langle u,z\rangle
   &=\tfrac{16}{3}r^2c+\tfrac{28}{3}cs,\\
 4\langle z,xz\rangle+\langle u,z\rangle
   &=\tfrac{16}{3}r^2c+\tfrac83cs.
 \end{aligned}
\]
Subtracting, and then using \(DQ(x)[y]/2-xQ(x)=0\), yields
\begin{equation}\label{eq:cubic-radial-polarization-uz}
 \langle u,z\rangle=\tfrac43cs,\qquad
 xu+4w=4sx+\tfrac83cy.
\end{equation}
Equation \eqref{eq:cubic-radial-polarization-eta} gives
\(D=2F^2-F(\nabla F)\). Differentiating and applying the second
identity in \eqref{eq:cubic-radial-polarization-uz}, we obtain
\begin{equation}\label{eq:cubic-intrinsic-D-gradient}
 \nabla D=\tfrac c3y-\tfrac18xu=\tfrac12(w-sx).
\end{equation}

Write \(b=\langle z,w\rangle\) and \(U=|u|^2\). Pairing
\eqref{eq:cubic-radial-polarization} with \(u\), and using
\(\langle z,xu\rangle=4s^2+32rc^2/9-4b\), gives
\[
 b=s^2+r\delta+\frac U{16}.
\]
The coefficient of \(\tau^3\) in \(P(x+\tau y)=0\) is
\[
 2U+4\langle z,xu\rangle+4b
 =\tfrac43\bigl(r\eta+6c\langle y,z\rangle+3s^2\bigr).
\]
Substitution of \eqref{eq:cubic-radial-polarization-eta} and
\(\langle Q,u\rangle=0\) reduces this to
\(b=s^2-U/4+8r\delta/3\). The two expressions for \(b\) imply
\begin{equation}\label{eq:cubic-radial-polarization-norms}
 U=\tfrac{16}{3}r\delta,\qquad
 b=s^2+\tfrac43r\delta.
\end{equation}

Differentiate \(\langle u,z\rangle=4cs/3\) in the direction \(u\):
\[
 4\langle w,xu\rangle+\langle u,yu\rangle+2|xu|^2
 =4s\eta+\tfrac{16}{3}c\langle z,u\rangle.
\]
Here \(P(y)=0\) gives \(\langle u,yu\rangle=4s\eta/3\), while
\eqref{eq:cubic-radial-polarization-uz} gives
\[
 \langle w,xu\rangle=4st+\tfrac{32}{9}c^2s-4|w|^2,
 \qquad |xu|^2=16rs^2-32st+16|w|^2.
\]
Substitution yields
\begin{equation}\label{eq:cubic-radial-polarization-w}
 |w|^2=\tfrac73st-\tfrac43rs^2.
\end{equation}
Finally, \(\nabla F=y/2\), \(\nabla S=z\), and
\eqref{eq:cubic-intrinsic-D-gradient} give
\[
 \begin{aligned}
 \Gamma(F,D)&=\tfrac14(\langle u,z\rangle-sc)=cs/12=2SF,\\
 \Gamma(S,D)&=\tfrac12(b-s^2)=2r\delta/3=8RD,\\
 \Gamma(D,D)&=\tfrac14(|w|^2-2st+rs^2)=s\delta/12=4SD.
 \end{aligned}
\]
\end{proof}

\begin{proposition}[Closure under a scalar quadratic trace form]
\label{prop:cubic-scalar-killing-closure}
Suppose, in addition to the hypotheses of
\cref{lem:cubic-radial-transports}, that \(F\) is harmonic and
\(\tr(L_x^2)=kR\) for a constant \(k\). Writing \(N=\dim V\), one has
\begin{equation}\label{eq:cubic-scalar-killing-laplacians}
 \Delta S=2kR,\qquad
 \Delta D=\left(2k-\frac{2N}{3}\right)S,
 \qquad \cA_R(F)=\R[R,F,S,D].
\end{equation}
In particular, this algebra has at most four generators.
\end{proposition}

\begin{proof}
Polarization gives \(\tr(L_xL_v)=k\langle x,v\rangle\).
Differentiating \eqref{eq:cubic-radial-polarization} and dividing
by four yields the operator identity
\[
 L_z+L_xL_y+L_yL_x+2L_x^3-2rL_x
 -\tfrac23c\Id-2x\otimes y-2y\otimes x=0,
\]
where \((a\otimes b)v=\langle b,v\rangle a\).
Taking traces, using harmonicity, gives
\[
 \tr L_x^3=\alpha c,\qquad \alpha=N/3+2-k.
\]
Its derivative in the direction \(y\) is
\(\tr(L_yL_x^2)=\alpha s\). Therefore
\[
 \begin{aligned}
 \operatorname{div}(yz)
 &=2\tr(L_zL_x)+\tr L_y^2+2\tr(L_yL_x^2)
   =(3k+2\alpha)s,\\
 \operatorname{div}(sx)&=(N+4)s.
 \end{aligned}
\]
Taking the divergence in \eqref{eq:cubic-intrinsic-D-gradient}
proves the formula for \(\Delta D\); Bochner's identity gives that
for \(\Delta S\). Together with \(\Delta R=2N\), \(\Delta F=0\),
homogeneity, and \cref{lem:cubic-radial-transports}, these formulas
show that \(\R[R,F,S,D]\) is Laplacian. Conversely,
\(S=\Gamma(F,F)\) and \eqref{eq:cubic-unified-D} put \(S,D\) in
every Laplacian algebra containing \(R,F\), proving equality.
\end{proof}

\begin{remark}[The radial bound and its trace input]
\label{rem:cubic-harmonic-generator-bound}
An irreducible harmonic radial eigencubic of exceptional type,
in the terminology of \cref{sec:cubic-history}, has \(\tr(L_x^2)=k|x|^2\)
for a constant \(k\), in the multiplication convention
\eqref{eq:cubic-polarized-product}. This is
\cite[Proposition~6.11.1]{NadirashviliTkachevVladut14}; see also
\cite[Proposition~5.2, equation~(5.11)]{TkachevClassification10}.
The scalar trace condition is not a consequence of radiality alone:
in the Clifford family, \eqref{eq:DeltaS} is scalar precisely when
\(m=q+1\).

Consequently every irreducible radial eigencubic has at most four
ordinary algebra generators in \(\cA_R(F)\). After normalizing
\(F\), the exceptional branch follows from this trace theorem and
\cref{prop:cubic-scalar-killing-closure}. In the Clifford branch,
the same proposition applies when \(m=q+1\); when
\(m\ne q+1\), \cref{thm:clifford-minimal-algebra} recovers the
block radii and identifies the smallest algebra directly.
\end{remark}


\subsubsection{The Clifford family and its differential table}
\label{sec:invariants}
We write the Clifford family of
\cite[Theorem~B]{PengXiao93} in the symmetric-system notation of
\cite{TkachevClifford10}.

Let \(q\geq1\), and let
\(\mathcal A=(A_0,\ldots,A_q)\) be a symmetric Clifford system on
\(\R^{2m}\), so that
\begin{equation}\label{eq:clifford-system}
 A_iA_j+A_jA_i=2\delta_{ij}\Id_{\R^{2m}}.
\end{equation}
For \((y,z)\in\R^{2m}\oplus\R^{q+1}\), put
\(A_z=\sum_{i=0}^qz_iA_i\).  The corresponding cubic and its zero cone are
\begin{equation}\label{eq:clifford-F}
 F=F_{\mathcal A}(y,z)=\sum_{i=0}^qz_i\langle A_i y,y\rangle
 =y^TA_zy,\qquad \mathcal C_{\mathcal A}=\{F=0\}.
\end{equation}
The Freudenthal--Springer algebra of \(F/2\) is a polar algebra in the
normalization of
\cite[Definition~6.5.1 and Propositions~6.5.1--6.5.2]{NadirashviliTkachevVladut14}.
We call the pair \((m,q)\) admissible when such a Clifford system exists.

Set
\[
 R=|y|^2+|z|^2,\qquad S=|\nabla F|^2,\qquad N=2m+q+1.
\]
The matrices are trace free, so \(F\) is harmonic.  The Clifford
relations give \(\cL(F)=-8RF\), and hence the regular zero set is
minimal; the calculations are given below.
We orient the cone as \(\partial\llbracket\{F<0\}\rrbracket\)
and use the local mass-minimizing convention for integral-current
competitors with compactly supported difference that bounds.

\paragraph{Differential identities.}

Introduce
\[
 \tau_i=y^TA_i y,\qquad
 \tau=(\tau_0,\ldots,\tau_q),\qquad
 A_\tau=\sum_{i=0}^{q}\tau_iA_i,
\]
and the scalar functions
\[
 u=|y|^2,\qquad v=|z|^2,\qquad R=u+v,\qquad T=|\tau|^2.
\]
The Clifford relations imply
\[
 A_z^2=v\Id,\qquad A_\tau^2=T\Id,\qquad
 A_zA_\tau+A_\tau A_z=2F\Id.
\]
Consequently,
\begin{equation}\label{eq:clifford-contractions}
 |A_zy|^2=uv,\quad |A_\tau y|^2=uT,\quad
 \langle A_zy,A_\tau y\rangle=uF.
\end{equation}

The vectors $A_0y,\ldots,A_qy$ are orthogonal, each of length $|y|$.
Bessel's inequality, followed by Cauchy--Schwarz, gives
\begin{equation}\label{eq:box-constraints}
 0\leq T\leq u^2,\qquad 0\leq F^2\leq vT.
\end{equation}
We shall also use
\[
 \langle A_i y,A_z y\rangle=uz_i,\qquad
 \langle A_i y,A_\tau y\rangle=u\tau_i.
\]

Set
\[
 U=\nabla F,\qquad S=|U|^2,\qquad Z=\nabla S,\qquad
 J=|Z|^2,\qquad K=\Delta S.
\]

The four functions \(I=(u,v,T,F)\) have the following gradient
inner products and Laplacians:
\begin{equation}\label{eq:clifford-closure-table}
 \begin{aligned}
 \bigl(\langle\nabla I_i,\nabla I_j\rangle\bigr)&=
 \begin{pmatrix}
 4u&0&8T&4F\\
 0&4v&0&2F\\
 8T&0&16uT&8uF\\
 4F&2F&8uF&4uv+T
 \end{pmatrix},\\
 \Delta I&=(4m,\,2(q+1),\,8(q+1)u,\,0).
 \end{aligned}
\end{equation}
For example, \(\nabla_yT=4A_\tau y\), so
\(|\nabla T|^2=16uT\) and \(\langle\nabla T,\nabla F\rangle=8uF\).
Euler's identity and the Clifford relations give the remaining entries.
Since \(R=u+v\), \cref{prop:MR-generator-test} shows that
\(\R[u,v,T,F]\) is a Laplacian algebra.
Hence \cref{thm:finite-reduction} applies to potentials \(G(u,v,T,F)\).

\begin{proposition}[Gradients and Laplacians]
\label{prop:first-order}
One has
\begin{subequations}
\label{eq:U}\label{eq:S}\label{eq:Z}\label{eq:Clifford-gradS-square}
\label{eq:DeltaF}\label{eq:DeltaS}\label{eq:UZ}
\begin{equation}
\begin{aligned}
 U&=(2A_z y,\tau),& S&=4uv+T,\\
 Z&=(8vy+4A_\tau y,8uz),& \Delta F&=0,\\
 K&=16mv+16(q+1)u,& \langle U,Z\rangle&=16RF,\\
 J&=16\bigl(4uvR+(u+4v)T\bigr).
\end{aligned}
\tag{\theparentequation}
\end{equation}
\end{subequations}
\end{proposition}

\begin{proof}
Differentiate to obtain $U,Z$; their squared norms follow from
\eqref{eq:clifford-contractions}. Conjugation by $A_j$, $j\ne i$,
sends $A_i$ to $-A_i$, proving trace-freeness and $\Delta F=0$.
The remaining formulas are
$K=\Delta(4uv+T)$ and $\langle U,Z\rangle=\Gamma(F,4uv+T)$
in \eqref{eq:clifford-closure-table}.
\end{proof}


\begin{theorem}[The smallest algebra of an unbalanced Clifford cubic]
\label{thm:clifford-minimal-algebra}
If \(m\ne q+1\), then
\[
 \cA_R(F)=\R[u,v,T,F],\qquad
 (\cA_R(F))_2=\operatorname{span}\{u,v\},\qquad
 (\cA_R(F))_3=\R F.
\]
The two coordinate blocks are therefore the intrinsic quadratic
blocks of \cref{thm:cubic-one-dimensional-degree-three}, with degree
pattern \((2,1)\).
\end{theorem}

\begin{proof}
Every Laplacian algebra containing \(R,F\) contains
\(S=\Gamma(F,F)\) and \(\Delta S\). Equation~\eqref{eq:DeltaS}
recovers the two block radii by
\[
 u=\frac{\Delta S/16-mR}{q+1-m},\qquad v=R-u.
\]
It then recovers \(T=S-4uv\). Thus \(\R[u,v,T,F]\) is
contained in \(\cA_R(F)\), and the reverse inclusion follows from
its closure table \eqref{eq:clifford-closure-table}. The generators
have degrees \(2,2,4,3\), giving the two homogeneous components
stated above; \(u,v\) are linearly independent.
Finally, if a constant vector \((\xi,\zeta)\) is a translation
direction of \(F\), its derivative has terms
\(2\langle A_z y,\xi\rangle\) and \(y^TA_\zeta y\) of
different bidegrees, so both vanish. Invertibility of \(A_0\)
gives \(\xi=0\), and \(A_\zeta^2=|\zeta|^2\Id\) gives
\(\zeta=0\). The block theorem therefore applies.
\end{proof}

When \(m=q+1\), one has $\Delta S=16mR$. Applying
\cref{prop:cubic-scalar-killing-closure} to the normalized cubic
$F/2$ therefore gives
\[
 \cA_R(F)=\R[R,F,S,vT]\subsetneq\R[u,v,T,F].
\]
Indeed, $|\nabla(F/2)|^2=S/4$, and its generator $D$ from
\eqref{eq:cubic-unified-D} is $vT/4$. In the present normalization,
\[
 vT=\frac{\Gamma(S,S)-16RS}{48}\in\cA_R(F).
\]
The generators $R,F,S,vT$ have degrees $2,3,4,6$, so the quadratic
component of this smallest algebra is $\R R$. Since $u,v$ are
linearly independent, the inclusion is strict. The larger algebra
remains useful for the coordinate calculations below.

Since \(Z=2\Hess F(U)\), \eqref{eq:UZ} and harmonicity give
\[
 \cL(F)=-\tfrac12\Gamma(F,S)=-8RF,
\]
recovering the radial minimality identity for this family
directly from the closure table.


\subsection{The weight operator and nonradial spectral rigidity}
\label{sec:nonradial-cubics}
\label{sec:cubic-generator-bound}

\begin{theorem}[The five-generator bound]
\label{thm:harmonic-cubic-five-bound}
Let $F$ be a nonzero irreducible homogeneous harmonic cubic on a
Euclidean space $V$, with $V=V_{\mathrm{ess}}(F)$, and suppose its
regular zero set is minimal.
Then the smallest Laplacian algebra $\cA_R(F)$ is
generated by at most five polynomials. More precisely:
\begin{enumerate}[label=\textup{(\roman*)}]
\item If $F$ is radial, $\cA_R(F)$ has at most four generators.
\item If $F$ is nonradial, $\cA_R(F)$ is a polynomial algebra on five
algebraically independent homogeneous generators of degrees
$(2,2,3,3,4)$. Its spherical quotient dimension is four.
\end{enumerate}
The upper bound five is attained by the examples of
\cref{sec:nonradial-models}.
\end{theorem}

Extending the Peirce framework of
\cite[Sections~6.6--6.8]{NadirashviliTkachevVladut14} to the weight
operator gives spectral rigidity and a graded splitting. Irreducibility
forces an exact eikonal core. The five resulting invariants close under
the differential operations, are recovered from $R,F$, and have rank
five. Spectral exclusion and core exactness are proved in
\cref{app:nr-spectral-exclusions,app:nr-core-exactness-proof}.

% Spectral structure used in the intrinsic five-generator bound.
\subsubsection{Maximal idempotents and the weight operator}
\label{subsec:nr-maximal-idempotents}

By \cref{lem:square-free-critical-set}, an irreducible harmonic $F$
has a nonsingular real zero. Thus \cref{prop:divisibility}
gives $\cL(F)=-qF$, where homogeneity makes $q$ quadratic.

Use the commutative Euclidean metrized algebra and normalization of
\cref{prop:cubic-multiplication-dictionary}.
Harmonicity is the condition $\tr L_x=0$ for every $x$;
we call such an algebra \emph{exact}, following
\cite[Section~3.2 and Definition~4.1]{FoxTkachev24}.
For a self-adjoint operator \(T\), the equation
\(\cL(F)=-\frac32\ip{Tx}{x}F\) is then equivalent to
\begin{equation}\label{eq:nr-Hsiang-identity}
 \ip{x^2}{x^3}=\ip{Tx}{x}\ip{x^2}{x}.
\end{equation}
This defines an \emph{exact Hsiang algebra}, with symmetric bilinear weight
$b(x,y)=\ip{Tx}{y}$. The radial case is the special case in which
$T$ is scalar.
Here \(x^3=x(x^2)\) and \(x^4=x(x^3)\); no associativity is assumed.
Differentiating \eqref{eq:nr-Hsiang-identity} gives
\begin{equation}\label{eq:nr-vector-identity}
 4x^4+(x^2)^2-2\ip{x^2}{x}Tx-3\ip{Tx}{x}x^2=0.
\end{equation}

As shown in \cref{subsec:cubic-multiplication},
$\operatorname{Ann}(A)=\{v:vA=0\}$ is the cylindrical subspace of $F$.
Metric invariance makes $A_{\mathrm{ess}}=\operatorname{Ann}(A)^\perp$
a subalgebra; its orthogonal splitting from the annihilator has zero
mixed products.

We use the spherical maximum construction of an idempotent
\cite[Lemma~6.6.1]{NadirashviliTkachevVladut14} and derive the spectral
bounds below.

\begin{lemma}[Maximal normalization]\label{lem:nr-max-idempotent}
Let \(A\) be a commutative metrized algebra with a Euclidean metric
and nonzero multiplication. After multiplying its product
by a positive constant, there is a unit idempotent \(c\) such that
\begin{equation}\label{eq:nr-spherical-maximum}
 \max_{|x|=1}\ip{x^2}{x}=\ip{c^2}{c}=1.
\end{equation}
Every eigenvalue of \(L_c\) on \(c^\perp\) lies in
\([-1,\frac12]\).
\end{lemma}

\begin{proof}
Put \(h(x)=\ip{x^2}{x}\). Polarization shows that \(h\) is a
nonzero odd cubic, so its maximum \(M\) on the unit sphere is positive.
At a maximizer \(c\), the Lagrange multiplier equation is \(c^2=Mc\).
Replacing the product by \(xy/M\) proves
\eqref{eq:nr-spherical-maximum}. It preserves exactness and replaces
\(T\) in \eqref{eq:nr-Hsiang-identity} by \(T/M^2\).
The constrained Hessian of \(h\) at \(c\) is
\(6L_c-3\Id\) on \(c^\perp\), giving the upper bound.
For a unit eigenvector \(u\perp c\), with \(cu=ru\), maximality gives
\[
 h\left(-\frac c2\pm\frac{\sqrt3}2u\right)
 =-\frac18-\frac98r\pm\frac{3\sqrt3}8h(u)\le1.
\]
Averaging gives \(r\ge-1\).
\end{proof}

Fix this normalization, which rescales the cubic by a positive constant.

\begin{lemma}[The weight operator at a unit idempotent]
\label{lem:nr-idempotent-quotient}
Every unit idempotent \(e\) in an algebra satisfying
\eqref{eq:nr-Hsiang-identity} satisfies
\begin{equation}\label{eq:nr-quotient-polynomial}
 Te=e,\qquad T=p(L_e)-6P_e,\qquad p(t)=4t^3+4t^2-t,
\end{equation}
where \(P_ex=\ip{e}{x}e\). In particular, \([T,L_e]=0\).
\end{lemma}

\begin{proof}
At \(e\), the scalar identity gives \(\ip{Te}{e}=1\), and the
vector identity gives \(Te=e\). Its derivative at \(e\) is
\[
 8L_e^3y+8L_e^2y-2L_ey-2Ty-12\ip{e}{y}e=0.
\]
This proves \eqref{eq:nr-quotient-polynomial}. Since \(L_e\) is
self-adjoint and \(L_ee=e\), it commutes with \(P_e\).
\end{proof}

Write \(E_r=\ker(L_c-r\Id)\). By maximality \(E_1=\R c\), and
\begin{equation}\label{eq:nr-Peirce-quotient}
 Tu=p(r)u\qquad (u\in E_r\subset c^\perp).
\end{equation}

\subsubsection{Fusion identities and spectral rigidity}
\label{subsec:nr-spectral-rigidity}

We adapt the radial polarization argument in
\cite[Lemma~6.6.2]{NadirashviliTkachevVladut14}, specifically equation~(6.6.8).
For nonradial $T$, we derive the fusion rule directly from
\eqref{eq:nr-vector-identity}, using
\begin{equation}\label{eq:nr-fusion-polynomial}
 \mathscr K(r,s,t)=4(r^2+s^2+t^2+rs+rt+st)+2(r+s+t)-3.
\end{equation}

\begin{lemma}[Fusion and the square identity]\label{lem:nr-fusion-square}
For \(u\in E_r\), \(v\in E_s\), \(w\in E_t\), all
perpendicular to \(c\), one has
\begin{equation}\label{eq:nr-fusion}
 \mathscr K(r,s,t)\ip{uv}{w}=0.
\end{equation}
For \(u\in E_r\subset c^\perp\), set
\(Q_r(t)=\mathscr K(r,r,t)\). Then
\begin{equation}\label{eq:nr-square-operator}
 Q_r(L_c)u^2=3\bigl(2r+p(r)\bigr)|u|^2c.
\end{equation}
Writing \(u^2=r|u|^2c+\sum_tq_t\), with
\(q_t\in E_t\cap c^\perp\), gives
\begin{equation}\label{eq:nr-square-norm}
 \sum_t(t+4r)|q_t|^2=4r^2(3r-1)(r+1)|u|^4.
\end{equation}
\end{lemma}

\begin{proof}
The mixed coefficient of \(ab\) in
\eqref{eq:nr-vector-identity} at \(c+au+bv\), projected to
\(c^\perp\) and divided by two, is the projection of
\[
 \bigl[4L_c^2+(4r+4s+2)L_c
       +4(r^2+rs+s^2)+2(r+s)-3\bigr]uv.
\]
Pairing with \(w\) proves \eqref{eq:nr-fusion}. The coefficient
of \(a^2\) at \(c+au\) gives
\eqref{eq:nr-square-operator}; thus \(Q_r(t)q_t=0\).
For the norm identity, pair the vector identity at \(u\) with \(c\).
Metric invariance gives
\[
 \ip{u^4}{c}=r|u^2|^2,\qquad
 \ip{(u^2)^2}{c}=r^2|u|^4+\sum_t t|q_t|^2.
\]
The term containing \(Tu\) pairs to zero, and the last term is
\(-3p(r)r|u|^4\). Rearranging yields
\eqref{eq:nr-square-norm}, since
\(3p(r)r-4r^3-r^2=4r^2(3r-1)(r+1)\).
\end{proof}

\begin{proposition}[A uniform spectral interval]
\label{prop:nr-spectral-interval}
Put \(r_0=-(2+\sqrt7)/6\). Then
\begin{equation}\label{eq:nr-Peirce-interval}
 \operatorname{Spec}(L_c|_{c^\perp})
 \subset\{-1,0\}\cup[r_0,-1/6]\cup[1/3,1/2].
\end{equation}
Moreover,
\begin{equation}\label{eq:nr-quotient-interval}
 \operatorname{Spec}T\subset\{0\}\cup
 \left[\frac7{27},\frac{17+7\sqrt7}{27}\right],
 \qquad \ker T=\operatorname{Ann}(A).
\end{equation}
\end{proposition}

\begin{proof}
All possible output eigenvalues \(t\) in
\eqref{eq:nr-square-norm} lie in \([-1,1/2]\). If
\(-1<r<r_0\), then \(Q_r(t)=4t^2+(8r+2)t+12r^2+4r-3\)
is decreasing on that interval and
\(Q_r(1/2)=12r^2+8r-1>0\). All \(q_t\) therefore vanish,
contradicting the negative right side of
\eqref{eq:nr-square-norm}. If \(-1/6<r<0\), both endpoint
values
\[
 Q_r(-1)=(6r+1)(2r-1),\qquad Q_r(1/2)=12r^2+8r-1
\]
are negative. Convexity makes \(Q_r\) negative throughout the
interval, giving the same contradiction. If \(0<r<1/3\),
then \(Q_r(-1)<0\) and \(Q_r(0)<0\); hence every admissible
root is positive. The left side of the norm identity is nonnegative
and the right side is negative. This proves
\eqref{eq:nr-Peirce-interval}. At \(r=1/3\), the norm identity
also shows that every \(q_t\) vanishes, because \(t+4/3>0\).

Applying \(p\) to these intervals gives the spectral inclusion in
\eqref{eq:nr-quotient-interval}. The only zero of \(p\) allowed
by \eqref{eq:nr-Peirce-interval} is \(0\). Thus \(u\in\ker T\)
implies \(cu=0\), and \eqref{eq:nr-square-operator} gives
\[
 (4L_c^2+2L_c-3\Id)u^2=0.
\]
Its two roots \((-1\pm\sqrt{13})/4\) lie outside
\([-1,1/2]\), and it is nonzero at \(1\). Consequently \(u^2=0\).
The coefficient of \(t^3\) in the vector identity at \(x+tu\)
is then \(8L_u^3x\): all other terms have degree at most two,
since \(Tu=u^2=0\). Hence \(L_u^3=0\), and self-adjointness
gives \(L_u=0\). Conversely, an annihilator element is orthogonal
to \(c=c^2\) and is killed by \(L_c\), so
\eqref{eq:nr-quotient-polynomial} kills it as well.
\end{proof}

By \cref{prop:nr-spectral-interval}, restriction to $A_{\mathrm{ess}}$
preserves exactness and the defining identity, and
$T|_{A_{\mathrm{ess}}}$ is positive definite.

Let
\begin{equation}\label{eq:nr-eigenvalue-one-space}
 B=E_T(1)=\R c\oplus Z\oplus H\oplus W,
 \qquad Z=E_{-1},\quad H=E_{-1/2},\quad W=E_{1/2},
\end{equation}
since \(p(t)-1=(t+1)(2t-1)(2t+1)\).
The fusion polynomial gives the following possible outputs on
\(c^\perp\):
\[
\begin{array}{c|cccccc}
(r,s)&(-1,-1)&(-1,-\frac12)&(-1,\frac12)&
(-\frac12,-\frac12)&(-\frac12,\frac12)&(\frac12,\frac12)\\ \hline
 t&\text{none}&\frac12&-\frac12,\frac12&-\frac12&
 -1,\frac12&-1,-\frac12
\end{array}
\]
(The additional root \(1\) in the fourth column is absent on
\(c^\perp\).) In particular,
\begin{equation}\label{eq:nr-B-subalgebra}
 B^2\subset B,\qquad BB^\perp\subset B^\perp.
\end{equation}

A Euclidean commutative metrized algebra is \emph{eikonal} if
\(x^3=|x|^2x\), equivalently \(|x^2|^2=|x|^4\).
Equivalently, the associated cubic $6F$ is a cubic eiconal in the sense
of \cite[Definition~3.1.1]{NadirashviliTkachevVladut14}; exactness is a
separate condition.

\begin{lemma}[The negative Peirce spaces]\label{lem:nr-negative-spaces}
\begin{enumerate}[label=\textup{(\roman*)}]
\item If \(H\ne0\), then \(\operatorname{Spec}T\subset\{1,0\}\).
\item If \(Z\ne0\), then
\(\operatorname{Spec}T\subset\{1,7/27,0\}\).
For \(U=E_T(7/27)\), one has
\begin{equation}\label{eq:nr-U-peirce}
 U=P\oplus Q,\quad P=E_{1/3},\quad Q=E_{-1/6},\qquad
 L_z|_P=0,\quad (L_z|_Q)^2=\frac1{12}|z|^2\Id_Q.
\end{equation}
\end{enumerate}
\end{lemma}

The proof is given in \cref{app:nr-spectral-exclusions}.

Exactness excludes $Z=H=W=0$.

\begin{lemma}\label{lem:nr-not-one-dimensional}
In an exact Hsiang algebra whose multiplication is not identically zero,
\(E_T(1)\ne\R c\).
\end{lemma}

The proof is given in \cref{app:nr-spectral-exclusions}.

\begin{theorem}[Spectral rigidity]\label{thm:nonradial-spectral}
Every exact Euclidean Hsiang algebra whose multiplication is not
identically zero admits a positive rescaling for which
\begin{equation}\label{eq:nr-rigid-spectrum}
 \operatorname{Spec}T\subset\{0,7/27,1\},\qquad
 \ker T=\operatorname{Ann}(A).
\end{equation}
On $A_{\mathrm{ess}}$, either $T|_{A_{\mathrm{ess}}}=\Id$ or
\begin{equation}\label{eq:nr-quotient-splitting}
 A_{\mathrm{ess}}=B\oplus U,\qquad T|_B=\Id,\qquad T|_U=\frac7{27}\Id.
\end{equation}
\end{theorem}

\begin{proof}
The kernel assertion is \cref{prop:nr-spectral-interval}.
If \(H\ne0\) or \(Z\ne0\), use
\cref{lem:nr-negative-spaces}.
Otherwise the fusion table gives \(w^2=|w|^2c/2\) for \(w\in W\).
The norm identity with \(r=1/2\) then reads
\(0=3|w|^4/4\), so \(W=0\). This would give
\(B=\R c\), contrary to \cref{lem:nr-not-one-dimensional}.
\end{proof}

% Intrinsic representation identities and exactness of the nonradial core.
\subsection{The eikonal core and its representation}
\label{subsec:nr-recovery}

Work on $A_{\mathrm{ess}}$, denoted by $A$, with nonradial weight operator.
By \cref{thm:nonradial-spectral} and \eqref{eq:nr-B-subalgebra},
\begin{equation}\label{eq:nr-recovery-splitting}
 A=B\oplus U,\qquad T|_B=\Id,\qquad T|_U=\tfrac7{27}\Id,
 \qquad B^2\subset B,\quad BU\subset U.
\end{equation}

\begin{proposition}[The eikonal core and its grading]\label{prop:nr-eikonal-core}
Assume $T|_{A_{\mathrm{ess}}}$ is not scalar.
Then \(H=0\), \(Z\ne0\), and the subalgebra \(B=E_T(1)\)
is eikonal. Moreover, $U^2\subset B$, so $A=B\oplus U$ is a
$\mathbb Z_2$-grading.
\end{proposition}

\begin{proof}
\Cref{lem:nr-negative-spaces}(i) and
\cref{thm:nonradial-spectral} give \(H=0\).
If \(Z=0\), the proof of \cref{thm:nonradial-spectral} would give
\(B=\R c\), so \(Z\ne0\). Fusion and invariance give
\[
 zz'=-\ip{z}{z'}c,\qquad ZW\subset W,\qquad
 w^2=\frac12|w|^2c+z_w,\quad z_w\in Z.
\]
For a unit \(z\in Z\), the self-adjoint map \(L_z\) preserves
\(W\). If \(w\in W\) is an eigenvector with eigenvalue \(y\),
then \(L_{e_\pm}w=(-1/4\pm\sqrt3y/2)w\), where
\(e_\pm=(-c\pm\sqrt3z)/2\). Since \(w\perp e_\pm\) and
\(Tw=w\), the quotient formula gives
\[
 p(-1/4+\sqrt3y/2)=p(-1/4-\sqrt3y/2)=1.
\]
The roots of \(p(t)-1\) are \(-1,-1/2,1/2\). Two such roots
with sum \(-1/2\) must be \(-1,1/2\), so \(y^2=3/4\).
Consequently \((L_z|_W)^2=3\Id/4\).
Scaling and the norm identity at \(w\in W\) yield
\[
 |zw|^2=\frac34|z|^2|w|^2,\qquad |z_w|^2=\frac34|w|^4.
\]
For \(b=tc+z+w\), its square has the orthogonal components
\[
 b^2=\left(t^2-|z|^2+\frac12|w|^2\right)c
          +(-2tz+z_w)+(tw+2zw).
\]
On taking norms, the terms linear in \(t\) cancel because
\(\ip{z}{z_w}=\ip{zw}{w}\). The displayed norm identities then
give \(|b^2|^2=|b|^4\). Differentiating within the metrized
subalgebra \(B\) proves \(b^3=|b|^2b\).

For the grading, use $U=P\oplus Q$ from \eqref{eq:nr-U-peirce}.
For the input eigenvalue pairs $(1/3,1/3)$, $(-1/6,-1/6)$,
and $(1/3,-1/6)$, the fusion rule \eqref{eq:nr-fusion}
permits, among the eigenvalues on $c^\perp$, respectively no output,
only $-1$, and only $1/2$.
Metric invariance therefore gives
\[
 p^2=\tfrac13|p|^2c,\qquad q^2\in\R c\oplus Z,\qquad pq\in W
 \qquad(p\in P,\ q\in Q).
\]
Polarization yields $U^2\subset B$; the other two grading relations
are already contained in \eqref{eq:nr-B-subalgebra}.
\end{proof}

At the maximal unit idempotent $c$, the decomposition is therefore
\[
 B=\R c\oplus Z\oplus W,\qquad U=P\oplus Q,
\]
with $L_c$-eigenvalues $1,-1,1/2,1/3,-1/6$, respectively, and $Z\ne0$.
To determine the action of $B$ on $U$, set $M_b=L_b|_U$ and
$\{A_1,A_2\}=A_1A_2+A_2A_1$.

\begin{proposition}\label{prop:nr-representation-identities}
For $b,d\in B$,
\begin{equation}\label{eq:nr-representation-quadratic}
\begin{aligned}
 M_{b^2}&=6M_b^2-\tfrac13|b|^2\Id_U,\\
 M_{bd}&=3\{M_b,M_d\}-\tfrac13\ip bd\Id_U.
\end{aligned}
\end{equation}
\begin{equation}\label{eq:nr-representation-cubic}
 108M_b^3-9|b|^2M_b-\ip{b^2}b\Id_U=0.
\end{equation}
\end{proposition}
\begin{proof}
Every unit idempotent $e\in B$ satisfies $\ip{e^2}e=1$, so it
attains the global spherical maximum in
\eqref{eq:nr-spherical-maximum}.
The maximality bounds and \cref{lem:nr-idempotent-quotient}, applied
on $U\subset e^\perp$, give
\[
 \operatorname{Spec}M_e\subset\{1/3,-1/6\}.
\]
Indeed, the third root $-7/6$ of $p(t)=7/27$ is below the lower
bound $-1$. Hence
$M_e=6M_e^2-\Id_U/3$.

Consider the homogeneous quadratic operator polynomial \(\mathscr Q(b)=M_{b^2}-6M_b^2+\tfrac13|b|^2\Id_U\).
Every nonzero $b\in B$ lies in an eikonal subalgebra of dimension
at most two. Indeed, differentiating $b^3=|b|^2b$ in direction $b^2$
gives
\[
 (b^2)^2=2\ip{b^2}b\,b-|b|^2b^2,
\]
so $\operatorname{span}\{b,b^2\}$ is such a subalgebra.
In dimension one it is generated by a unit idempotent, and hence
$\mathscr Q$ vanishes there.
In dimension two choose a positive spherical maximum $e$ for its cubic.
Since $|e^2|=1$, the multiplier equation gives $e^2=e$ without
rescaling. For a unit vector $z\perp e$, write
$ez=\alpha z$, $z^2=\alpha e+\beta z$.
Linearization at $e$ gives $2\alpha^2+\alpha=1$, while $z^3=z$
gives $\alpha\beta=0$ and $\alpha^2+\beta^2=1$.
Thus $\alpha=-1$, $\beta=0$, and
\[
 e,\qquad \frac{-e+\sqrt3z}{2},\qquad
 \frac{-e-\sqrt3z}{2}
\]
are three pairwise nonproportional unit idempotents.
A homogeneous quadratic polynomial on a plane vanishing on three
distinct lines is zero. Applying this to each matrix entry of
$\mathscr Q$ proves $\mathscr Q(b)=0$.
Polarization gives the second quadratic identity.
Finally, apply it to $b,d=b^2$. Since $b^3=|b|^2b$,
\[
 |b|^2M_b
 =3\{M_b,6M_b^2-\tfrac13|b|^2\Id_U\}
       -\tfrac13\ip{b^2}b\Id_U
 =36M_b^3-2|b|^2M_b-\tfrac13\ip{b^2}b\Id_U,
\]
which gives the cubic identity.
\end{proof}

The remaining identities on $U$ are as follows.

\begin{lemma}\label{lem:nr-U-cubic}
For $u\in U$, set $s=|u|^2$, $w=u^2\in B$, and $v=wu\in U$.
Then
\begin{equation}\label{eq:nr-U-identities}
 v=\tfrac s9u,\qquad |w|^2=\tfrac{s^2}9,\qquad
 w^2=\tfrac s3w,\qquad uv=\tfrac s9w.
\end{equation}
Moreover, for every $b\in B$,
\begin{equation}\label{eq:nr-mixed-polarization}
 -3wb+18u(bu)=sb.
\end{equation}
\end{lemma}
\begin{proof}
Pair the first identity in \eqref{eq:nr-representation-quadratic}
with $u$ and differentiate in $b$ to obtain
\eqref{eq:nr-mixed-polarization}.
At $u$, \eqref{eq:nr-vector-identity} reads
$4uv+w^2=7sw/9$.
Together with \eqref{eq:nr-mixed-polarization} at $b=w$, this gives
$uv=sw/9$ and $w^2=sw/3$.
Since $B$ is eikonal, $|w^2|^2=|w|^4$, and therefore
\[
 |u^2|^2\left(|u^2|^2-\tfrac19|u|^4\right)=0
 \qquad(u\in U).
\]
The first factor is not the zero polynomial: $U\ne0$ and the
only possible eigenvalues $1/3,-1/6$ of $M_c$ are nonzero, so
$\ip{u^2}c=\ip{M_cu}u$ is not identically zero.
The polynomial ring on $U$ is an integral domain, whence
$|u^2|^2=|u|^4/9$ identically.
Differentiating this identity within $U$ gives $4wu=4su/9$.
\end{proof}

\begin{proposition}[Exactness of the core]\label{prop:nr-core-exactness}
If the original cubic is irreducible over $\R$, then $B$ is exact,
\[
 \tr_B L_b=\tr_U M_b=0\qquad(b\in B),
\]
and $\dim B=\dim U+2$.
\end{proposition}
The proof is given in \cref{app:nr-core-exactness-proof}.

\subsection{The intrinsic five-generator closure}
\label{subsec:nr-five-generators}

Use the normalization in \cref{lem:nr-max-idempotent}.

\begin{theorem}[Five generators in the nonradial branch]
\label{thm:nonradial-five-generators}
Let $F$ be a nonzero irreducible homogeneous harmonic cubic on a
Euclidean space $V=V_{\mathrm{ess}}(F)$, with minimal regular zero set
and nonradial quadratic weight. After a constant rescaling of $F$, let
$V=B\oplus U$ be the orthogonal splitting of
\cref{thm:nonradial-spectral}. For $x=b+u$, define
\begin{equation}\label{eq:nr-five-invariants}
 r=|b|^2,\qquad s=|u|^2,\qquad
 a=\ip{b^2}b,\qquad p=\ip{bu}u,\qquad h=|bu|^2.
\end{equation}
Then
\begin{equation}\label{eq:nr-five-algebra}
 \cA_R(F)=\R[r,s,a,p,h],\qquad
 R=r+s,\qquad F=\frac{a+3p}{6}.
\end{equation}
These five polynomials are algebraically independent. In particular,
the smallest Laplacian algebra of $F$ requires exactly five ordinary
algebra generators, of degrees $2,2,3,3,4$.
\end{theorem}

\begin{proof}
By \cref{thm:nonradial-spectral,prop:nr-eikonal-core},
both $B$ and $U$ are nonzero, $B$ is eikonal, and
$B^2\subset B$, $BU\subset U$, $U^2\subset B$.
By \cref{prop:nr-core-exactness}, irreducibility makes $B$ exact,
and the symmetric operators $M_b=L_b|_U$ satisfy
\begin{equation}\label{eq:nr-five-traces}
 \tr_B L_b=0,\qquad \tr_U M_b=0\qquad(b\in B).
\end{equation}
Write $m=\dim U$, $n=\dim B=m+2$, and $w=u^2\in B$.
The identities of
\cref{prop:nr-representation-identities,lem:nr-U-cubic}, together with
\eqref{eq:nr-mixed-polarization}, give
\begin{align}
 M_{b^2}&=6M_b^2-\frac r3\Id_U,
 &108M_b^3&=9rM_b+a\Id_U,\label{eq:nr-five-representation}\\
 |w|^2&=\frac{s^2}{9},
 &w^2&=\frac s3w,\notag\\
 M_wu&=\frac s9u,
 &18u(bu)&=3bw+sb.\label{eq:nr-five-mixed}
\end{align}
The eikonal identity on $B$ gives $b(b^2)=rb$ and $|b^2|^2=r^2$.

\emph{Differential closure.}
The five bidegrees with respect to $B\oplus U$ are, respectively,
\[
 (2,0),\quad(0,2),\quad(3,0),\quad(1,2),\quad(2,2).
\]
Thus, for any of the
five generators $f$ of bidegree $(i,j)$,
\[
 \Gamma(r,f)=2if,\qquad \Gamma(s,f)=2jf.
\]
All remaining gradient products are
\begin{equation}\label{eq:nr-five-gamma}
\begin{aligned}
 \Gamma(a,a)&=9r^2,\\
 \Gamma(a,p)&=18h-rs,\\
 \Gamma(p,p)&=\frac{s^2}{9}+4h,\\
 \Gamma(a,h)&=rp+\frac{sa}{3},\\
 \Gamma(p,h)&=\frac{rp}{3}+\frac{2sp}{9}+\frac{sa}{27},\\
 \Gamma(h,h)&=\frac{r+s}{3}h+\frac{p^2}{9}+\frac{ap}{27}.
\end{aligned}
\end{equation}
To verify the table intrinsically, metric invariance and
\eqref{eq:nr-five-mixed} give the gradients, written in their
$B\oplus U$ components,
\begin{equation}\label{eq:nr-five-gradients}
 \nabla a=(3b^2,0),\qquad
 \nabla p=(w,2M_bu),\qquad
 \nabla h=\left(\frac{bw}{3}+\frac{sb}{9},\,2M_b^2u\right).
\end{equation}
Their pairings, using \eqref{eq:nr-five-representation}, yield the
first five lines of \eqref{eq:nr-five-gamma}.
For the last line, polarize the eikonal identity on $B$ at $w$ to obtain
\[
 2w(wb)+w^2b=|w|^2b+2\ip wb w.
\]
Pairing with $b$ and using
$\ip{b^2}w=6h-rs/3$ gives
\[
 |bw|^2=\frac{s^2r}{9}+p^2-sh.
\]
Consequently
\[
 |\nabla_Bh|^2=\frac{sh}{3}+\frac{p^2}{9},\qquad
 |\nabla_Uh|^2=4\ip{M_b^4u}u
             =\frac{rh}{3}+\frac{ap}{27},
\]
which proves the final pairing.

The Laplacians are
\begin{equation}\label{eq:nr-five-laplacians}
 \Delta r=2(m+2),\qquad \Delta s=2m,\qquad
 \Delta a=\Delta p=0,\qquad
 \Delta h=\frac{mr+(m+2)s}{9}.
\end{equation}
To verify these formulas, exactness of $B$ implies
$\sum_{\alpha=1}^n e_\alpha^2=0$ for an orthonormal basis of $B$.
Summing the quadratic identity in \eqref{eq:nr-five-representation},
and separately taking its trace, gives
\[
 \sum_{\alpha=1}^n M_{e_\alpha}^2=\frac n{18}\Id_U,
 \qquad \tr M_b^2=\frac m{18}r.
\]
Consequently,
\[
 \Delta_Bh=2\sum_{\alpha=1}^n|M_{e_\alpha}u|^2=\frac n9s,
 \qquad \Delta_Uh=2\tr M_b^2=\frac m9r.
\]
The other formulas in \eqref{eq:nr-five-laplacians} follow from
\eqref{eq:nr-five-traces}. The differential chain rule now proves
that $\R[r,s,a,p,h]$ is a Laplacian algebra.

\emph{Recovery from $R$ and $F$.}
The grading gives $F=(a+3p)/6$. Applying
\eqref{eq:nr-five-gamma} and \eqref{eq:nr-five-laplacians} gives
\begin{equation}\label{eq:nr-five-gradient-square}
 S=\frac14(r-s/3)^2+4h,\qquad
 \Delta S=\frac{(10m+36)r+(2m+4)s}{9}.
\end{equation}
Since $R=r+s$ and $m>0$, we recover the two radii directly:
\begin{equation}\label{eq:nr-five-radius-recovery}
 r=\frac{9\Delta S-(2m+4)R}{8(m+4)},\qquad s=R-r.
\end{equation}
Both belong to $\cA_R(F)$, because $S=\Gamma(F,F)$ and
$\Delta S$ do. Setting $D=\Gamma(r,F)=a+p$ then gives
\[
 p=3F-\frac D2,\qquad
 a=\frac{3D}{2}-3F,\qquad
 h=\frac{\Gamma(a,p)+rs}{18}.
\]
Thus all five generators belong to $\cA_R(F)$. Since their algebra
is Laplacian and contains $R,F$, equality follows.

\emph{Algebraic independence.}
At the maximal unit idempotent $c$,
\cref{prop:nr-eikonal-core} supplies a nonzero $(-1)$-Peirce space
$Z\subset B$. Choose a unit vector $z\in Z$. Then $z^2=-c$, so at
$b=z$ one has $r=1$, $a=0$, and
\[
 \nabla_B r=2z,\qquad \nabla_B a=-3c.
\]
In particular, the differential of $(r,a)$ has rank two.
By \eqref{eq:nr-five-representation}, every eigenvalue of $M_z$ lies in
\[
 \{\lambda_1,\lambda_2,\lambda_3\}
   =\left\{-\frac1{2\sqrt3},0,\frac1{2\sqrt3}\right\}.
\]
Let their multiplicities be $k_1,k_2,k_3$, allowing zero provisionally.
The identities $\tr M_z=0$, $\tr M_z^2=m/18$, and
$k_1+k_2+k_3=m$ imply $k_1=k_2=k_3=m/3>0$.
Choose $u=u_1+u_2+u_3$ with a nonzero component in each eigenspace.
For fixed $b=z$, the map from $t_i=|u_i|^2$ to the remaining invariants
is
\[
 \begin{pmatrix}s\\p\\h\end{pmatrix}
 =
 \begin{pmatrix}
  1&1&1\\
  \lambda_1&\lambda_2&\lambda_3\\
  \lambda_1^2&\lambda_2^2&\lambda_3^2
 \end{pmatrix}
 \begin{pmatrix}t_1\\t_2\\t_3\end{pmatrix}.
\]
The Vandermonde matrix is invertible, and the differentials $dt_i$
are independent because $u_i\ne0$. Thus $(s,p,h)$ has rank three
in the $U$ directions, whereas $(r,a)$ depends only on $b$ and has
rank two. The full five-invariant map has rank five at this point.
Its components are therefore algebraically independent. Their
polynomial algebra has transcendence degree five and cannot be
generated by fewer than five elements.
\end{proof}

\begin{corollary}[Cartan reduction and ambient dimensions]
\label{cor:nr-cartan-dimensions}
Under the hypotheses and normalization of
\cref{thm:nonradial-five-generators}, the cubic
\begin{equation}\label{eq:nr-cartan-reduction}
 H=a-9p=6\bigl(\Gamma(r,F)-5F\bigr)
\end{equation}
satisfies
\[
 \Delta H=0,\qquad |\nabla H|^2=9R^2.
\]
It is orthogonally equivalent to a normalized Cartan cubic, and
\[
 \dim V\in\{8,14,26\}.
\]
Writing $C=H/6$ and $C_B=C\circ\operatorname{pr}_B$, one has
\begin{equation}\label{eq:nr-cartan-reconstruction}
 F=\frac{4C_B-C}{3}.
\end{equation}
\end{corollary}

\begin{proof}
The formula for $H$ follows from $6F=a+3p$ and
$\Gamma(r,F)=a+p$. By \eqref{eq:nr-five-gradients},
\[
 \nabla H=(3b^2-9u^2,-18bu).
\]
The identities $|b^2|^2=r^2$, $|u^2|^2=s^2/9$, and
$\ip{b^2}{u^2}=6h-rs/3$ therefore give
\[
 |\nabla H|^2
 =9r^2+9s^2-54(6h-rs/3)+324h=9(r+s)^2.
\]
Also $\Delta H=0$ by \eqref{eq:nr-five-laplacians}.

Put $m=\dim U$. The multiplicity calculation in the proof of
\cref{thm:nonradial-five-generators} gives $m/3\in\mathbb N$
and $m>0$. Since $\dim B=m+2$, it follows that
$\dim V=2m+2\ge8$.
Cartan's classification of harmonic cubic eikonals
\cite{Cartan39}, or \cite[Theorem~1.1]{TkachevCartan10},
now identifies $H$ with a Cartan cubic in dimension $5,8,14$, or
$26$. In the latter theorem, the additional reducible possibility
$x_N^3-3x_N\sum_{i<N}x_i^2$ has Laplacian $6(2-N)x_N$,
so harmonicity and $N\ge8$ exclude it. The same dimension bound
excludes $N=5$. Finally, $C_B=a/6$ and $C=(a-9p)/6$ give
\eqref{eq:nr-cartan-reconstruction}.
\end{proof}

The reduction recovers $F$ from a Cartan cubic and a compatible
splitting $B\oplus U$. Classifying $F$ up to equivalence therefore
also requires classifying these splittings.

The models in \cref{sec:nonradial-models} realize the sharp bound
and express \(r,s\) and the other invariants in coordinates.
The programs check their closure, recovery, and rank-five minor;
the universal spectral and trace arguments are proved above and in
\cref{app:cubic-structure}.


\begin{proof}[Proof of \cref{thm:harmonic-cubic-five-bound}]
The divisibility argument in \cref{subsec:nr-maximal-idempotents}
and the cubic correspondence associate an exact Hsiang algebra to $F$.
In the radial case, write $\cL(F)=-\gamma RF$, where $\gamma>0$
by the maximal normalization above. Harmonicity gives
$\Gamma(F,S)=2\gamma RF$. Replacing $F$ by
$\widetilde F=\sqrt{2/\gamma}\,F$ gives
$\Gamma(\widetilde F,|\nabla\widetilde F|^2)=4R\widetilde F$,
as required by \cref{prop:cubic-scalar-killing-closure}; this rescaling
does not change $\cA_R(F)$.
The branch argument in \cref{rem:cubic-harmonic-generator-bound}
then gives the radial four-generator bound.

In the nonradial case, \cref{thm:nonradial-five-generators} gives the
polynomial algebra with the stated generator degrees. Restriction to
$R=1$ reduces its quotient dimension from five to four. The examples
in \cref{prop:nonradial-algebra} establish sharpness.
\end{proof}
